\section{Implementierung}

\subsection{Knowledge}

\subsubsection{Motivation und Zielsetzung}

Der Knowledge Graph dient der semantischen Modellierung
physiotherapeutischer Übungsdaten. Ziel ist es, Übungen
maschinenlesbar mit Körperteilen, Kategorien,
Schwierigkeitsgraden und Kontraindikationen zu verknüpfen,
um eine intelligente Empfehlungslogik zu ermöglichen.

\subsubsection{Ontologie-Design}

\paragraph{Verwendete Ontologie}

Als Basis wurde die \textit{Physical Therapy Ontologie}
(\texttt{http://purl.org/PhysicalTherapy}) verwendet und
durch projektspezifische Vokabulare erweitert.

\subsection{API und Datenhaltung}

\subsection{UI und Interaktion}

\subsection{Integration}
